% docs/manual/developer/chapters/02-dev-environment.tex
% 第 2 章：开发环境

\chapter{开发环境}

本章给从 \emph{零安装} 到 \emph{能跑测试 + 能改代码} 的最短路径。读完后你应当能：clone 仓库、安装可编辑模式、跑 ruff + pytest、用 superpowers 流程提交修改。

\section{Fork 与 clone}

\openbench{} 主仓库在 \href{https://github.com/CoLM-SYSU/OpenBench}{github.com/CoLM-SYSU/OpenBench}。贡献流程是 fork → branch → PR：

\begin{minted}{bash}
# 1. 在 GitHub 上点 Fork
# 2. clone 你的 fork
git clone https://github.com/<your-username>/OpenBench.git
cd OpenBench
# 3. 加 upstream 便于同步
git remote add upstream https://github.com/CoLM-SYSU/OpenBench.git
git remote -v
\end{minted}

之后 sync 上游：

\begin{minted}{bash}
git fetch upstream
git merge upstream/main
\end{minted}

\section{Python 与虚拟环境}

\openbench{} 要求 Python 3.10+。强烈建议用虚拟环境隔离：

\subsection{方式 1：uv（推荐，最快）}

\begin{minted}{bash}
# 装 uv
curl -LsSf https://astral.sh/uv/install.sh | sh

# 在项目根目录建 venv
uv venv .venv --python 3.12
source .venv/bin/activate

# 可编辑安装 + 全部 extras + 开发工具
uv pip install -e ".[all,dev]"
\end{minted}

\subsection{方式 2：conda}

\begin{minted}{bash}
conda create -n openbench python=3.12
conda activate openbench
pip install -e ".[all,dev]"
\end{minted}

\subsection{方式 3：标准 venv}

\begin{minted}{bash}
python3 -m venv .venv
source .venv/bin/activate
pip install -e ".[all,dev]"
\end{minted}

\section{各 extras 含义}

\file{pyproject.toml} 定义了 4 个可选 extras：

\begin{itemize}
  \item \cli{[gui]}：PySide6（向导 GUI）
  \item \cli{[remote]}：paramiko + cryptography（SSH 远程执行）
  \item \cli{[report]}：xhtml2pdf（PDF 报告生成）
  \item \cli{[dev]}：pytest、pytest-cov、ruff
  \item \cli{[all]}：合并 \cli{[gui,remote,report]}，但 \emph{不含} \cli{[dev]}
\end{itemize}

开发者完整安装命令：

\begin{minted}{bash}
pip install -e ".[all,dev]"
\end{minted}

\warninline{\cli{[gui]} 在无图形系统的 HPC 节点装不上 PySide6。如果你只在服务器上开发非 GUI 部分，用 \cli{".[remote,report,dev]"}。}

\section{验证环境}

\begin{minted}{bash}
# 版本
openbench version
# 预期：openbench 3.0.0a1

# CLI 入口
openbench --help

# 跑一次完整测试套件
pytest tests/ -v

# 编译手册（可选）
cd docs/manual && make probe && make all
\end{minted}

\section{Ruff lint + format}

\openbench{} 用 \modname{ruff} 同时做 lint 与 format。配置在 \file{pyproject.toml} 的 \texttt{[tool.ruff]} 段。

\subsection{运行}

\begin{minted}{bash}
# Lint（仅检查）
ruff check src tests

# Format（自动修）
ruff format src tests

# 一并修可修的 lint 问题
ruff check --fix src tests
\end{minted}

\subsection{规则集}

\begin{itemize}
  \item \texttt{E}：pycodestyle 错误
  \item \texttt{F}：pyflakes（未使用 import 等）
  \item \texttt{I}：isort（import 排序）
  \item \texttt{W}：pycodestyle 警告
\end{itemize}

\subsection{已迁移文件的豁免}

\file{pyproject.toml} 中 \texttt{[tool.ruff.lint.per-file-ignores]} 列出了 v1.x 迁移过来的"老代码"豁免：

\begin{minted}{toml}
"src/openbench/util/*.py" = ["E", "F", "W", "I"]
"src/openbench/core/*.py" = ["E", "F", "W", "I"]
# ...
\end{minted}

\textbf{含义}：这些路径下的代码暂时不强制 lint。原因是 v3.0 是 v1.x 的 Python 化迁移，强行 lint 会引发大量 churn。新代码（\file{config/}、\file{runner/}、\file{cli/}、\file{tests/}）仍然必须过 lint。

\noteinline{开新文件、改新代码时，务必让 ruff 干净。修老文件时只在涉及行附近修 lint 问题，不做大规模重排。}

\section{Pytest 与 pytest-cov}

\subsection{测试目录}

\begin{itemize}
  \item \file{tests/}：pyproject 里 \yamlkey{testpaths} 指定的目录
  \item \file{tests/test\_config/}、\file{tests/test\_registry/}、\file{tests/test\_runner/}、\file{tests/manual/}
\end{itemize}

\subsection{常用命令}

\begin{minted}{bash}
# 跑全部测试
pytest

# 详细输出 + 跑某文件
pytest tests/test_config/test_loader.py -v

# 按模式选
pytest -k "init"   # 只跑名字含 init 的测试

# 看覆盖率
pytest --cov=openbench --cov-report=term-missing

# 失败立刻停
pytest -x

# 并发跑（要装 pytest-xdist）
pytest -n auto
\end{minted}

\subsection{pythonpath}

\file{pyproject.toml}:

\begin{minted}{toml}
[tool.pytest.ini_options]
testpaths = ["tests"]
pythonpath = ["src", "."]
\end{minted}

\texttt{src} 让 \texttt{import openbench} 找到 \file{src/openbench/}；\texttt{.} 让 \texttt{import docs.manual.scripts.*}（手册生成器）也能 import。

\section{Pre-commit hooks}

如果项目根有 \file{.pre-commit-config.yaml}（可选）：

\begin{minted}{bash}
pip install pre-commit
pre-commit install
# 之后每次 git commit 自动跑 ruff + 其他 hook
\end{minted}

如果没有，手动在 commit 前跑 ruff：

\begin{minted}{bash}
ruff check src tests && ruff format --check src tests && pytest
\end{minted}

\section{Superpowers 工作流}

\openbench{} 用 \modname{superpowers} 系列文档化中大型改动，分四个阶段：

\begin{enumerate}
  \item \textbf{Spec}（\file{docs/superpowers/specs/}）：设计文档；说"为什么改、改什么、不改什么"
  \item \textbf{Plan}（\file{docs/superpowers/plans/}）：实施计划；按 task 拆，每 task 有可运行代码与测试
  \item \textbf{Execute}：按 plan 落地；每个 task 单独 commit
  \item \textbf{Review}（\file{docs/superpowers/reviews/}）：审计与回顾；记录发现的 bug 与决议
\end{enumerate}

\subsection{何时用，何时不用}

\textbf{用}：

\begin{itemize}
  \item 新功能（新 CLI 命令、新 metric、新 GUI page）
  \item 架构改动（schema 改动、API 重命名）
  \item 修复需要细化讨论的 bug
  \item 文档项目（如本手册）
\end{itemize}

\textbf{不用}：

\begin{itemize}
  \item 单 commit 能描述清的小 bug（直接 PR）
  \item Typo / docstring 修复
  \item Test 新增（已有功能加测试）
\end{itemize}

\subsection{文件命名}

\begin{itemize}
  \item Spec：\file{YYYY-MM-DD-<topic>-design.md}
  \item Plan：\file{YYYY-MM-DD-<topic>.md}
  \item Review：\file{YYYY-MM-DD-<topic>.md}（同名）
\end{itemize}

详见第 11 章贡献流程。

\section{编辑器与 IDE}

\openbench{} 不强制特定 IDE，但推荐配置：

\subsection{VS Code}

推荐扩展：

\begin{itemize}
  \item Python（Microsoft）
  \item Pylance（类型检查）
  \item Ruff（实时 lint + format）
  \item GitLens（提交历史可视化）
\end{itemize}

\file{.vscode/settings.json} 建议：

\begin{minted}{json}
{
  "python.testing.pytestEnabled": true,
  "python.testing.pytestArgs": ["tests"],
  "[python]": {
    "editor.formatOnSave": true,
    "editor.defaultFormatter": "charliermarsh.ruff"
  }
}
\end{minted}

\subsection{PyCharm}

\begin{itemize}
  \item Settings → Project → Interpreter → 选 \file{.venv/bin/python}
  \item Tools → External Tools → 加 ruff
  \item Run → Edit Configurations → 加 pytest
\end{itemize}

\section{常见开发环境问题}

\subsection{cartopy 装不上}

第 1 章用户卷已详解。简言之：

\begin{minted}{bash}
brew install geos proj   # macOS
sudo apt install libgeos-dev libproj-dev   # Linux
\end{minted}

或用 conda 装 cartopy：

\begin{minted}{bash}
conda install -c conda-forge cartopy
pip install -e ".[remote,report,dev]"   # 跳过 [all] 中的 cartopy
\end{minted}

\subsection{import openbench 找不到}

通常是 \texttt{pip install -e .} 没装成功，或者忘了 activate venv。检查：

\begin{minted}{bash}
which python   # 确认指向 .venv
pip list | grep -i openbench   # 确认装了
\end{minted}

\subsection{测试导入失败}

通常是 \texttt{pythonpath} 没配好。看 \file{pyproject.toml [tool.pytest.ini\_options]}：

\begin{minted}{toml}
pythonpath = ["src", "."]
\end{minted}

如果你 \texttt{cd tests/} 跑 pytest，要回到仓库根：\texttt{cd ..}。

\section{下一步}

下一章深入 \modname{config} 子系统：schema dataclass、loader 流程、adapter 翻译、migration 路径、resolver 解析。这是 \openbench{} 从用户输入到执行引擎之间最复杂的一层。
